% =====================================================================
\chapter{Metodologia}\label{cap:metodologia}
% =====================================================================

Este capítulo descreve o delineamento da pesquisa, a amostra e as fontes de dados,
o critério de classificação das empresas familiares, os indicadores de
endividamento e os procedimentos estatísticos empregados para avaliar a hipótese
H1.

% ---------------------------------------------------------------------
\section{Abordagem e estratégia da pesquisa}\label{sec:abordagem}
% ---------------------------------------------------------------------

A pesquisa adota abordagem quantitativa, de natureza descritiva e comparativa, com
uso de dados secundários de companhias abertas. O objetivo é verificar se empresas
de controle familiar e não familiar diferem quanto ao nível e à composição do
endividamento, e não estimar uma relação causal entre a estrutura de controle e a
estrutura de capital.

A proposta inicial previa um modelo de painel dinâmico estimado por System-GMM,
aplicado a todas as companhias não financeiras da B3. Ao longo da coleta, a
classificação manual das empresas a partir dos Formulários de Referência levou à
delimitação de um único setor (\autoref{sec:amostra}), o que resultou em uma
amostra de 16 companhias. Com esse número de unidades, estimadores de painel
dinâmico baseados em instrumentos internos não são adequados, pois exigem um
número de empresas consideravelmente maior que o de instrumentos. Optou-se, por
isso, pela comparação entre grupos por meio de estatística descritiva e de testes
de diferença entre grupos, complementados por análises de subperíodos e de
robustez.

% ---------------------------------------------------------------------
\section{Definição da amostra e fontes de dados}\label{sec:amostra}
% ---------------------------------------------------------------------

A população do estudo compreende as companhias abertas classificadas pela B3 no
setor econômico de Saúde. A restrição a um único setor tem duas finalidades:
reduzir a influência de diferenças setoriais na estrutura de capital, que a
literatura aponta como determinante relevante do endividamento
\cite{oliveira2020}, e viabilizar a classificação manual de cada companhia quanto
ao controle familiar, que exige a leitura individual dos Formulários de
Referência.
\verificar{incluir outro motivo para a escolha do setor de saúde, se houver}

Os dados contábeis foram extraídos da base Economática, em valores consolidados,
com frequência trimestral, de 31/12/2018 a 31/12/2025 (29 datas-base). Os
indicadores de fluxo, como EBITDA e despesa financeira, correspondem aos últimos
12 meses de cada data-base. Das 22 companhias do setor com dados na base, foram
mantidas apenas as que possuem informações em todo o intervalo, para que os dois
grupos sejam comparados no mesmo período. Foram excluídas seis empresas que
abriram capital ou passaram a ter dados durante o período: D1000 (DMVF, dados a
partir de 12/2019), Eurofarma (EUFA), Oncoclínicas (ONCO) e Viveo (VVEO), a partir
de 12/2020, Mater Dei (MATD), a partir de 09/2021, e Pague Menos (PGMN), a partir
de 09/2022. A amostra final é composta por 16 companhias, distribuídas em quatro
subsetores da B3 (\autoref{qua:amostra}).

\begin{quadro}[htb]
\caption{Composição da amostra por subsetor e classificação (Cenário 3)}\label{qua:amostra}
\centering\small
\begin{tabularx}{\textwidth}{|P{4.6cm}|L|L|}
\hline
\textbf{Subsetor (B3)} & \textbf{Familiares} & \textbf{Não familiares} \\ \hline
Comércio e Distribuição & Dimed/Panvel (PNVL), Profarma (PFRM), RaiaDrogasil (RADL) & -- \\ \hline
Medicamentos e Outros Produtos & Blau (BLAU), Hypera (HYPE), Ourofino (OFSA) & Biomm (BIOM) \\ \hline
Serviços Médico-Hospitalares, Análises e Diagnósticos & Dasa (DASA), Hapvida (HAPV), Rede D'Or (RDOR) & Alliar (AALR), Bradsaúde (SAUD), Fleury (FLRY), Qualicorp (QUAL) \\ \hline
Equipamentos & Baumer (BALM) & Lifemed (LMED) \\ \hline
\textbf{Total} & \textbf{10} & \textbf{6} \\ \hline
\end{tabularx}
\fonte{Elaborado pelo autor (2026), com base na classificação setorial da B3 e nos
Formulários de Referência (CVM).}
\end{quadro}

A série da Bradsaúde (SAUD) corresponde, no período analisado, à antiga Odontoprev
(ODPV3), cuja denominação e código de negociação foram alterados. A companhia não
apresenta dívida financeira em nenhuma das datas-base e foi mantida na amostra
principal, com análise de robustez que a exclui (\autoref{sec:procedimentos}).
\verificar{confirmar na DFP a ausência de empréstimos e financiamentos da SAUD e a
dívida bruta nula da HAPV em 31/12/2018}

% ---------------------------------------------------------------------
\section{Classificação das empresas familiares}\label{sec:classificacao}
% ---------------------------------------------------------------------

Não há definição única de empresa familiar na literatura, e os estudos brasileiros
combinam, de formas diferentes, critérios de propriedade, de controle e de
participação da família na administração. Para esta pesquisa, foram reunidos os
critérios de seis trabalhos (\autoref{qua:estudos}), e a partir deles foram
construídos três cenários de classificação, que diferem principalmente no
percentual mínimo de ações ordinárias detido pela família.

\begin{quadro}[htb]
\caption{Critérios de classificação de empresa familiar na literatura}\label{qua:estudos}
\centering\small
\begin{tabularx}{\textwidth}{|P{3.4cm}|L|L|P{3.2cm}|}
\hline
\textbf{Estudo} & \textbf{Participação mínima} & \textbf{Envolvimento na administração} & \textbf{Amostra} \\ \hline
\citeonline{santos2018} & Mais de 15\% das ações ordinárias, com a família no controle & Familiar no conselho de administração & 110 empresas da B3, 2010--2016 \\ \hline
\citeonline{tommasetti2019} & Mais de 50\% dos votos & Familiar como diretor-presidente ou presidente do conselho & 151 empresas da B3, 2012--2016 \\ \hline
\citeonline{voelcker2020} & Família como acionista controladora declarada & Membro indicado pelo controlador ao conselho & 1.188 observações da B3, 2010--2017 \\ \hline
\citeonline{poletti2022} & Pelo menos 20\% das ações & Familiar conselheiro, diretor-presidente ou presidente do conselho & 187 empresas da América Latina, 2005--2016 \\ \hline
\citeonline{ribeiro2026} & Pelo menos 10\% das ações ordinárias & Familiar na administração (dimensões tratadas separadamente) & 105 empresas da B3, 2010--2024 \\ \hline
\citeonline{ey2024} & Pelo menos 32\% dos votos (companhia aberta) & Não exige & 500 maiores empresas familiares do mundo \\ \hline
\end{tabularx}
\fonte{Elaborado pelo autor (2026), com base nos estudos citados.}
\end{quadro}

Em todos os cenários, a empresa é classificada como familiar somente se atender,
simultaneamente, a três condições: (i) a família detém um percentual mínimo das
ações ordinárias, somando a participação direta e a indireta; (ii) a família
integra o controle declarado da companhia; e (iii) há membro da família na
administração. A empresa que deixa de atender a pelo menos uma condição é
classificada como não familiar. O \autoref{qua:cenarios} apresenta os três
cenários.

\begin{quadro}[htb]
\caption{Cenários de classificação adotados}\label{qua:cenarios}
\centering\small
\begin{tabularx}{\textwidth}{|P{2.2cm}|P{2.2cm}|L|L|P{2.3cm}|}
\hline
\textbf{Cenário} & \textbf{Ações ordinárias da família} & \textbf{Administração} & \textbf{Base} & \textbf{Familiares $\times$ não familiares} \\ \hline
Cenário 1 & Mais de 50\% & Familiar como diretor-presidente ou presidente do conselho & \citeonline{tommasetti2019} & 5 $\times$ 11 \\ \hline
Cenário 2 & Pelo menos 32\% & Familiar no conselho de administração & \citeonline{ey2024}; \citeonline{voelcker2020} & 7 $\times$ 9 \\ \hline
\textbf{Cenário 3 (principal)} & \textbf{Mais de 15\%} & \textbf{Familiar no conselho de administração} & \textbf{\citeonline{santos2018}} & \textbf{10 $\times$ 6} \\ \hline
\end{tabularx}
\fonte{Elaborado pelo autor (2026). Em todos os cenários, exige-se também que a
família integre o controle declarado da companhia.}
\end{quadro}

O Cenário 3, que reproduz o critério de \citeonline{santos2018}, foi adotado como
classificação principal por três razões de método: trata-se de critério de estudo
brasileiro que utiliza a mesma fonte de informação desta pesquisa (Formulário de
Referência) e a mesma base de participação (ações ordinárias); capta o controle
exercido por blocos familiares com participação inferior a 50\%, frequentemente
organizados por acordo de acionistas, situação comum entre as companhias da B3; e,
de forma secundária, produz grupos de tamanho mais equilibrado, o que é relevante em uma amostra
pequena. Os Cenários 1 e 2 são utilizados como testes de robustez, para verificar
se os resultados se mantêm quando se exige uma participação familiar maior.

Para aplicar as condições, adotaram-se as seguintes regras, comuns aos três
cenários:

\begin{alineas}
\item \textbf{base de participação:} considera-se o percentual das ações
ordinárias, isto é, das ações com direito a voto, e não do capital total;

\item \textbf{grupo familiar:} a ``família'' pode corresponder a uma ou mais
famílias que integram o mesmo bloco de controle, inclusive por acordo de
acionistas, nos termos do art.~116 da Lei n.º 6.404/1976 \cite{brasil1976}. É o
caso da RaiaDrogasil (RADL), controlada por um bloco formado por três famílias
que, em conjunto, detêm cerca de 28\% das ações ordinárias e têm membros no
conselho de administração, o que a classifica como familiar no Cenário 3, mas não
nos Cenários 1 e 2;

\item \textbf{participação indireta:} soma-se a participação detida por meio de
\textit{holdings}, fundos de investimento em participações ou outras pessoas
jurídicas controladas pela família, seguindo a cadeia societária até as pessoas
físicas. Na Dasa (DASA), por exemplo, a participação da família Bueno, de cerca de
57\% das ações ordinárias, é detida por meio de um fundo de investimento em
participações;

\item \textbf{família no controle:} a participação relevante não basta se a
família não integra o controle declarado. Na Fleury (FLRY), a família Pardini
detém cerca de 20\% das ações ordinárias desde a incorporação do Instituto Hermes
Pardini, mas não faz parte do bloco de controle, e por isso a companhia é
classificada como não familiar em todos os cenários. Também são não familiares as
companhias de capital disperso ou controladas por instituições, como a Bradsaúde
(SAUD), controlada pelo Banco Bradesco;

\item \textbf{fonte e data de referência:} as informações foram obtidas no
Formulário de Referência mais recente disponível no site da Comissão de Valores
Mobiliários (CVM), nos itens de controle e composição acionária (itens 15.1 e 15.2
até 2022 e itens 6.1 e 6.2 a partir de 2023) e de administração, complementadas
pelos sites de relações com investidores. A classificação é mantida fixa em todo o
período, como em \citeonline{tommasetti2019}. Mudanças relevantes de controle ao
longo do período, como as reorganizações de Hapvida, Rede D'Or e Fleury em
2021--2022, são reconhecidas como limitação.
\end{alineas}

\verificar{percentuais de RADL (28,3\%), DASA (56,7\%) e FLRY (19,73\%) no
Formulário de Referência mais recente}

% ---------------------------------------------------------------------
\section{Indicadores de endividamento}\label{sec:indicadores}
% ---------------------------------------------------------------------

Na ausência de dados de valor de mercado e de um modelo de regressão, o
endividamento foi medido por indicadores contábeis calculados pela Economática,
organizados em três níveis de prioridade conforme a sua relação com a hipótese
(\autoref{qua:indicadores}).

\begin{quadro}[htb]
\caption{Indicadores de endividamento}\label{qua:indicadores}
\centering\small
\begin{tabularx}{\textwidth}{|P{2.5cm}|P{2.9cm}|L|L|}
\hline
\textbf{Prioridade} & \textbf{Indicador} & \textbf{Cálculo} & \textbf{O que mede} \\ \hline
Principal & Dívida bruta / Ativo total & (Empréstimos, financiamentos e debêntures de curto e longo prazo) $\div$ Ativo total & Endividamento oneroso: parcela do ativo financiada por credores financeiros, que impõem \textit{covenants} e monitoramento \\ \hline
Principal & Exigível total / Ativo total & (Passivo circulante + Passivo não circulante) $\div$ Ativo total & Endividamento geral, incluindo obrigações operacionais (fornecedores, salários, tributos, provisões) \\ \hline
Complementar & Dívida bruta / (Dívida bruta + PL) & Dívida bruta $\div$ (Dívida bruta + Patrimônio líquido) & Estrutura de capital: proporção de dívida no capital investido \\ \hline
Complementar & Dívida líquida / EBITDA & (Dívida bruta $-$ Caixa e aplicações) $\div$ EBITDA dos últimos 12 meses & Capacidade de pagamento; métrica usual em \textit{covenants} financeiros \\ \hline
Complementar & Dívida de curto prazo / Dívida total & Dívida onerosa de curto prazo $\div$ Dívida onerosa total & Composição da dívida por prazo e exposição ao risco de refinanciamento \\ \hline
Secundário & Exigível total / PL; Dívida bruta / PL; Dívida líquida / PL & Passivo ou dívida $\div$ Patrimônio líquido & Grau de endividamento em relação ao capital próprio \\ \hline
\end{tabularx}
\fonte{Elaborado pelo autor (2026), com base nas definições da Economática.}
\end{quadro}

A Dívida bruta / Ativo total é o indicador mais diretamente ligado à hipótese, pois
mede a dívida financeira, que traz os credores e as cláusulas restritivas
discutidos na \autoref{sec:risco}. O Exigível total / Ativo total permite verificar
se uma eventual diferença se restringe à dívida financeira ou alcança o passivo
como um todo. A composição da dívida por prazo foi incluída porque a dívida de
curto prazo exige renovações frequentes com os credores e, portanto, aumenta a
exposição da empresa à sua interferência \cite{diamond1991}.

Alguns indicadores não podem ser calculados em determinadas situações, e esses
casos foram tratados da seguinte forma: a Dívida de curto prazo / Dívida total não
é definida para a Bradsaúde, que não possui dívida financeira, e essa empresa é
excluída apenas desse indicador; a Dívida líquida / EBITDA não é informada pela
base quando o EBITDA é negativo, o que ocorre com a Biomm de 2018 a 2024. Nesse
indicador, a média da Biomm corresponde apenas ao exercício de 2025, e os resultados
são apresentados com e sem a empresa, acompanhados das medianas dos grupos, que são
menos sensíveis a valores extremos.
Os indicadores calculados sobre o patrimônio líquido (PL) são tratados como
secundários, por se tornarem instáveis quando o PL é pequeno ou negativo, como
ocorre com a Alliar em 2025.

% ---------------------------------------------------------------------
\section{Procedimentos de análise}\label{sec:procedimentos}
% ---------------------------------------------------------------------

A análise foi realizada em cinco etapas. Os cálculos foram feitos no Google
Planilhas e repetidos, de forma independente, em um \textit{script} em Python que
lê a mesma base e gera as tabelas deste trabalho; os dois procedimentos produzem
os mesmos resultados. A planilha de apoio, com a base de dados, a classificação
das empresas e as fórmulas, e o \textit{script} acompanham este trabalho.
\verificar{forma de disponibilização da planilha à banca: anexo, apêndice ou link}

\textbf{a) Análise descritiva.} Para cada indicador, calculou-se a média de cada
grupo em cada data-base, como média simples entre as empresas, de modo que todas as
companhias tenham o mesmo peso. A evolução dos grupos é apresentada em gráficos com
as 29 datas trimestrais, e as tabelas registram os valores de 31 de dezembro de
cada exercício, bem como o número de exercícios em que a média das familiares
ficou abaixo da média das não familiares.

\textbf{b) Testes de diferença entre os grupos.} A unidade de análise dos testes é
a empresa: para cada companhia, calculou-se a média do indicador nas oito datas de
encerramento de exercício (31/12/2018 a 31/12/2025), obtendo-se uma observação por
empresa. Esse procedimento foi adotado porque as observações de uma mesma empresa
em trimestres ou anos consecutivos não são independentes entre si, e tratá-las
como observações distintas superestimaria a significância das diferenças. Sobre
essas médias, aplicaram-se dois testes bicaudais:

\begin{itemize}
\item o teste \textit{t} de Welch \cite{welch1947}, que compara as médias dos dois
grupos sem exigir variâncias iguais nem grupos de mesmo tamanho; e
\item o teste U de Mann-Whitney \cite{mann1947}, não paramétrico, que compara as
posições (postos) das empresas de cada grupo e é menos sensível a valores
extremos e à ausência de normalidade, comuns em indicadores contábeis de amostras
pequenas. O p-valor foi obtido pela aproximação normal com correção de
continuidade; o p-valor exato, calculado para conferência, não altera nenhuma das
conclusões.
\end{itemize}

Adotaram-se como referência os níveis de significância de 1\%, 5\% e 10\%, e um
p-valor entre 5\% e 10\% é interpretado como evidência fraca. Embora a hipótese H1
seja direcional, os testes bicaudais foram mantidos por serem mais conservadores e
por a literatura apresentar argumentos nos dois sentidos (\autoref{sec:risco}).
Dado o número reduzido de empresas, que limita o poder dos testes, os resultados
não significativos são interpretados como ausência de evidência estatística de
diferença, e não como evidência de que os grupos são iguais
\cite{wasserstein2016}.

\textbf{c) Análise por subperíodos.} Para examinar se a diferença entre os grupos
se altera diante de um choque externo, o período foi dividido em três
subperíodos: pré-pandemia (2018--2019), pandemia de COVID-19 (2020--2021) e
pós-pandemia (2022--2025). A pandemia exigiu decisões de financiamento rápidas e
fora do planejamento das empresas, o que a torna um período de interesse para
observar a relação das empresas com os credores. Os testes da etapa (b) foram
repetidos com a média de cada empresa em cada subperíodo. Essa divisão foi definida
durante a análise, depois de observados os dados, e tem, por isso, caráter
exploratório: os resultados dos três subperíodos são reportados integralmente,
inclusive os não significativos, e, para o conjunto desses testes, informa-se o
p-valor ajustado pelo método de Holm.

\textbf{d) Análises de robustez.} Os testes foram repetidos (i) com as
classificações dos Cenários 1 e 2, para verificar a sensibilidade dos resultados ao
critério de empresa familiar; (ii) excluindo a Bradsaúde, operadora de planos
controlada por uma instituição financeira e única companhia da amostra sem dívida
financeira, para avaliar a influência dessa observação sobre o resultado, sem que
isso signifique tratar a ausência de dívida como erro, já que ela é uma escolha
legítima de financiamento; (iii) retirando uma empresa de cada vez, entre as 16,
para verificar se algum resultado depende de uma única companhia; e (iv) excluindo
as três empresas do subsetor de Comércio e Distribuição, em que não há empresas
não familiares.

\textbf{e) Magnitude das diferenças.} Além dos p-valores, são apresentados a
diferença entre as médias dos grupos, o seu intervalo de confiança de 95\% e o
tamanho de efeito (g de Hedges), para indicar a magnitude e a imprecisão das
estimativas.

Como foram realizados diversos testes (vários indicadores, dois testes,
subperíodos e cenários), um resultado significativo isolado deve ser interpretado
com cautela, e as conclusões consideram o conjunto dos resultados, e não testes
individuais.

% [OPCIONAL – só incluir se a regressão for refeita com o subsetor corrigido e um
% script reprodutível (R):]
% \textbf{e) Análise exploratória multivariada.} Como complemento, estimou-se uma
% regressão linear com dados empilhados (\textit{pooled}), tendo como variável
% dependente a Dívida bruta / Ativo total e como variável de interesse uma variável
% binária para empresa familiar, controlando por tamanho (logaritmo do capital
% empregado), rentabilidade do ativo, liquidez corrente, imobilização do patrimônio
% líquido e efeitos de ano, com erros-padrão agrupados por empresa. Dado o número
% de empresas, os resultados dessa etapa são tratados como exploratórios.
